\documentclass[10pt,a4paper]{amsart}
\usepackage[margin=19mm]{geometry}
\usepackage{amsmath,amssymb,mathtools}
\usepackage[scr=rsfs,cal=euler]{mathalfa}
\usepackage{xfrac}
\usepackage[unicode,hidelinks,bookmarks=false]{hyperref}
\newcommand{\ie}{i.e.\ }
\newcommand{\cf}{cf.\ }
\newcommand{\resp}{respectively\ }
\newcommand{\Subsection}{\subsection}
\providecommand{\nobreakdash}{\nobreak\hbox{-}}
\setlength{\emergencystretch}{2em}
\multlinegap0pt
\usepackage{amssymb}
\usepackage{xcolor}
\newtheorem{theorem}{Theorem}
\title{Unlocking Novel Topological Structures via Rough Families}
\author{Sourav Mandal, Lakshmi Kanta Dey, Pratikshan Mondal}
\date{Source: arXiv:2607.22018v1 (CC BY 4.0)}
\begin{document}
\maketitle

\noindent\emph{Source and licence.} Unlocking Novel Topological Structures via Rough Families. Version arXiv:2607.22018v1 (CC BY 4.0), distributed by arXiv under the Creative Commons Attribution 4.0 International licence, http://creativecommons.org/licenses/by/4.0/, as recorded in the arXiv licence metadata for this exact version and shown on the abstract page https://arxiv.org/abs/2607.22018v1. Full licence notice: \texttt{licenses/arxiv-2607.22018-CC-BY-4.0.txt}.\par\smallskip

\noindent\textit{Editorial scope.} Counted unit: the article's theorem theorem (source0.tex lines 162--173). The article's complete original proof of it is delivered with it (source0.tex lines 175--191, one \texttt{proof} environment of 17 lines). No mathematics is written, repaired or re-derived: the statement and the proof are the article's own lines, in the article's own notation, and no retained line is substituted or re-flowed.  No further source-local context is needed: every cross-reference of every retained line resolves inside the two delivered spans.  Nothing was omitted from any retained span, so the package carries no omission record.  No retained line contains a citation, so the wrapper carries no bibliography and no bibliography entry.  No retained line contains a figure environment, an image-inclusion command or drawing code, so no image is delivered and the package declares no asset dependency.  Editorial formatting only, in addition: an AMS \texttt{amsart} preamble that restates the article's own declarations and macros used by the retained lines, namely the environments \texttt{theorem} on separate counters, together with the standard packages those lines need.  The retained environments print in this document as Theorem 1 in order of appearance, whereas the article prints the same items with the numbering of its own full document.  The complete native source of the article as distributed in the arXiv e-print for this exact version is bundled unchanged as \texttt{sources/205906/source0.tex}.
\begin{theorem}
    Let $\mathscr{F}$ be a rough family in a topological space $(X,\tau)$. Then, for $A, ~B\subset X,$ the following statements are true:
    \begin{itemize}
        \item[(i)] If $A\subset B,$ then $\overline{A}^\mathscr{F}\subseteq \overline{B}^\mathscr{F},$ ~$A^{o\mathscr{F}}\subseteq B^{o\mathscr{F}}.$
        
        \item[(ii)] $\overline{A}^\mathscr{F}\cup\overline{B}^\mathscr{F}= \overline{(A\cup B)}^\mathscr{F}$ and $A^{o\mathscr{F}}\cup B^{o\mathscr{F}}\subset (A\cup B)^{o\mathscr{F}}.$ In the second case, equality might not happen.
        
        \item[(iii)] $\overline{(A\cap B)}^\mathscr{F}\subset \overline{A}^\mathscr{F}\cap\overline{B}^\mathscr{F}$ and $(A\cap B)^{o\mathscr{F}}=A^{o\mathscr{F}}\cap B^{o\mathscr{F}}.$ Usually, the equality might not be true for the first case.
        
        \item[(iv)] $\left(\overline{A}^\mathscr{F}\right)^c=(A^c)^{o\mathscr{F}},$ $\left(A^{o\mathscr{F}}\right)^c=\overline{(A^c)}^\mathscr{F}.$
    \end{itemize}
\end{theorem}

\begin{proof}
 \begin{itemize}
     \item[(i)]The proof follows directly from the definitions.
     
     \item[(ii)] The first part of (ii) follows from the definitions of rough closure jointly with (i). 
     
     For the subsequent part, let us consider $\mathbb{R}$ endowed with the lower limit topology $\mathbb{R}_l.$ Let $A=[0,1)$ and $B=[1,2).$ Then for the rough family $\mathscr{F}=\left\{\left\{x,2x\right\}: x\in \mathbb{R}\right\},$ we obtain $$A^{o\mathscr{F}}\cup B^{o\mathscr{F}}=\left[0,\frac{1}{2}\right)\cup\emptyset=\left[0,\frac{1}{2}\right)\neq \left[0,1\right)=(A\cup B)^{o\mathscr{F}}.$$
     
     \item[(iii)]  The first part of (iii) is also a consequence of (i) and the definition of rough interior.
     
     For the later part, we substitute the rough family $\mathscr{F}$ with $\mathscr{F}=\left\{\left\{x,\frac{x}{2}\right\}: x\in \mathbb{R}\right\}$ in the aforementioned scenario, yielding $$\overline{(A\cap B)}^\mathscr{F}=\overline{\emptyset}^\mathscr{F}=\emptyset\neq [1,2)=[0,2)\cap [1,4)=\overline{A}^\mathscr{F}\cap\overline{B}^\mathscr{F}.$$
     
     \item[(iv)]  For the first part, note that $x\in (A^c)^{o\mathscr{F}}$ if and only if there is $U\in \tau$ with $F_x\subset U\subset A^c,$ i.e., $(U\cap A)\setminus\{x\}=\emptyset$ which implies that $x\notin \overline{A}^\mathscr{F}$ or $x\in \left(\overline{A}^\mathscr{F}\right)^c$.

     For the second part, let $x\in \left(A^{o\mathscr{F}}\right)^c.$ Then, there exists $U\in \tau$ with $F_x\subset U$ but $\left(U\cap A^c\right)\setminus\{x\}\neq \emptyset$. Thus, $x\in \overline{(A^c)}^\mathscr{F}$ implies $\left(A^{o\mathscr{F}}\right)^c\subset \overline{(A^c)}^\mathscr{F}$. Similarly, we can obtain $\overline{(A^c)}^\mathscr{F}\subset \left(A^{o\mathscr{F}}\right)^c$ and we are done.
 \end{itemize}
\end{proof}

\end{document}
